\chapter[1973 --- Frisch et al.]{1973: Karl von Frisch, Konrad Lorenz, Nikolaas Tinbergen}
\label{ch:1973_Frisch}

\section*{Main Contribution}

\textit{Elucidated the organization of animal behavior, showing that even complex behaviors have innate components that can be studied scientifically.}

\section{Official Citation}

for their discoveries concerning organization and elicitation of individual and group behaviour

\section{Background and Historical Context}

The 1973 Nobel Prize in Physiology or Medicine was awarded to Karl von Frisch, Konrad Lorenz, and Nikolaas Tinbergen for their discoveries concerning organization and elicitation of individual and social behaviour patterns. This prize recognized the establishment of ethology as a scientific discipline---the study of animal behaviour in natural environments.

\subsection{The Rise of Ethology}

By the mid-20th century, the study of animal behaviour had become a major field of biological research. However, there were two main approaches to studying animal behaviour:

\begin{itemize}
\item Comparative psychology: This approach, dominant in the United States, focused on learning and conditioning and used laboratory experiments to study animal behaviour
\item Ethology: This approach, pioneered by von Frisch, Lorenz, and Tinbergen, focused on innate behaviour patterns and studied animals in their natural environments
\end{itemize}

The work of von Frisch, Lorenz, and Tinbergen established ethology as a rigorous scientific discipline and demonstrated the importance of studying animal behaviour in natural settings.

\subsection{Karl von Frisch}

Karl von Frisch was born on November 20, 1886, in Vienna, Austria. He studied at the University of Vienna, where he earned his doctorate in 1910. After completing his doctorate, von Frisch worked at the University of Rostock, the University of Breslau, and the University of Munich, where he would spend most of his career.

Von Frisch was a skilled zoologist who is best known for his discovery of the honey bee dance language---a sophisticated communication system by which honey bees convey information about the location of food sources to their nestmates.

\subsection{Konrad Lorenz}

Konrad Lorenz was born on November 7, 1903, in Vienna, Austria. He studied at the University of Vienna, where he earned his doctorate in 1933. After completing his doctorate, Lorenz worked at various research institutions in Austria and Germany, and later at the Max Planck Institute for Behavioural Physiology in Seewiesen, Germany, where he would spend most of his career.

Lorenz was a charismatic and controversial figure who is best known for his work on imprinting---a type of learning that occurs during a critical period early in life.

\subsection{Nikolaas Tinbergen}

Nikolaas Tinbergen was born on April 15, 1907, in The Hague, Netherlands. He studied at the University of Leiden, where he earned his doctorate in 1932. After completing his doctorate, Tinbergen worked at the University of Leiden and later at the University of Oxford, where he would spend most of his career.

Tinbergen was a skilled experimentalist who developed rigorous methods for studying animal behaviour. He is best known for his four questions, a framework for studying animal behaviour that considers causation, development, function, and evolution.

\section{The Discovery/Contribution}

The work of von Frisch, Lorenz, and Tinbergen on animal behaviour involved the discovery and analysis of innate behaviour patterns and the establishment of ethology as a scientific discipline.

\subsection{Von Frisch's Work on the Honey Bee Dance Language}

Karl von Frisch's primary contribution was the discovery of the honey bee dance language.

\subsubsection{The Discovery}

Von Frisch discovered that honey bees communicate the location of food sources to their nestmates through a sophisticated dance language. The dance consists of two main components:

\begin{itemize}
\item The waggle dance: A figure-eight dance in which the bee waggles its abdomen while moving in a straight line. The angle of the straight line relative to the vertical indicates the direction of the food source relative to the sun, and the duration of the waggle run indicates the distance to the food source.
\item The round dance: A circular dance that indicates that food is nearby but does not provide specific directional information.
\end{itemize}

\subsubsection{Significance of the Discovery}

Von Frisch's discovery of the honey bee dance language was significant because it demonstrated that animals can communicate complex information through symbolic means. The dance language is one of the most sophisticated communication systems known in the animal kingdom and has been the subject of extensive research since its discovery.

\subsection{Lorenz's Work on Imprinting}

Konrad Lorenz's primary contribution was the discovery and analysis of imprinting.

\subsubsection{The Discovery}

Lorenz discovered that certain types of learning occur during a critical period early in life and are irreversible. He called this type of learning imprinting. The most famous example of imprinting is the attachment of newly hatched goslings to their mother. If the goslings are exposed to Lorenz during the critical period, they will follow him instead of their biological mother.

\subsubsection{Key Features of Imprinting}

Lorenz identified several key features of imprinting:

\begin{itemize}
\item It occurs during a critical period early in life
\item It is irreversible---once imprinting has occurred, it cannot be changed
\item It is species-specific---imprinting occurs only to members of the same species
\item It involves both recognition and attachment
\end{itemize}

\subsubsection{Significance of the Discovery}

Lorenz's work on imprinting was significant because it demonstrated the importance of critical periods in development and provided insights into the mechanisms of learning and behaviour. The concept of critical periods has been applied to many areas of biology, including the study of human development.

\subsection{Tinbergen's Four Questions}

Nikolaas Tinbergen's primary contribution was the development of a framework for studying animal behaviour, known as Tinbergen's four questions.

\subsubsection{The Four Questions}

Tinbergen proposed that animal behaviour should be studied at four levels:

\begin{itemize}
\item Causation: What mechanisms (neural, hormonal, muscular) cause the behaviour?
\item Development: How does the behaviour develop during the lifetime of the individual?
\item Function: What is the adaptive value of the behaviour? How does it contribute to survival and reproduction?
\item Evolution: How has the behaviour evolved over evolutionary time?
\end{itemize}

\subsubsection{Significance of the Framework}

Tinbergen's four questions provided a comprehensive framework for studying animal behaviour that is still widely used today. This framework ensures that all aspects of behaviour are considered, from the underlying mechanisms to the evolutionary history.

\section{Impact on Modern Medicine}

The work of von Frisch, Lorenz, and Tinbergen on animal behaviour has had a significant impact on modern medicine, particularly in the fields of psychiatry, psychology, and neuroscience.

\subsection{Understanding Human Behaviour}

The study of animal behaviour has provided important insights into human behaviour. Many of the principles discovered by von Frisch, Lorenz, and Tinbergen apply to human behaviour as well as animal behaviour. For example, the concept of critical periods has been applied to the study of human development, including language acquisition and social development.

\subsection{Psychiatric Disorders}

The understanding of innate behaviour patterns has also contributed to our understanding of psychiatric disorders. Many psychiatric disorders, such as autism spectrum disorder and attention deficit hyperactivity disorder (ADHD), are thought to involve abnormalities in innate behaviour patterns. The principles established by von Frisch, Lorenz, and Tinbergen have been essential for understanding these conditions.

\subsection{Behavioural Therapy}

The understanding of animal behaviour has also contributed to the development of behavioural therapy for psychiatric disorders. Behavioural therapy uses principles of learning and behaviour modification to treat conditions such as phobias, anxiety disorders, and obsessive-compulsive disorder. The principles established by von Frisch, Lorenz, and Tinbergen have been essential for the development of these therapies.

\subsection{Animal Models of Disease}

The study of animal behaviour has also contributed to the development of animal models of disease. By studying behaviour in animals, researchers can gain insights into the mechanisms of human diseases and develop new treatments. The principles established by von Frisch, Lorenz, and Tinbergen have been essential for the development of these models.

\subsection{Conservation Medicine}

The understanding of animal behaviour has also contributed to conservation medicine---the application of medical knowledge to the conservation of wildlife. By understanding the behaviour of endangered species, researchers can develop strategies for their conservation. The principles established by von Frisch, Lorenz, and Tinbergen have been essential for this work.

\section{Legacy}

The legacy of Karl von Frisch, Konrad Lorenz, and Nikolaas Tinbergen extends far beyond their Nobel Prize-winning work on animal behaviour. Their discoveries established ethology as a scientific discipline and have contributed to numerous advances in medicine, psychology, and conservation biology.

\subsection{Foundation for Ethology}

The work of von Frisch, Lorenz, and Tinbergen established ethology as a rigorous scientific discipline. Their discoveries about innate behaviour patterns, imprinting, and the analysis of animal behaviour provided the tools and principles for subsequent research in ethology. The principles they established continue to guide research in animal behaviour today.

\subsection{Advances in Medicine}

The understanding of animal behaviour has contributed to numerous advances in medicine, including the development of behavioural therapies for psychiatric disorders, the understanding of human development, and the development of animal models of disease. These advances have improved the lives of millions of patients and continue to drive progress in these fields.

\subsection{Evolutionary Medicine}

The ethological approach to studying behaviour has also contributed to the field of evolutionary medicine---the application of evolutionary principles to understanding disease. By studying the evolutionary history of behaviour, researchers can gain insights into the mechanisms of disease and develop new treatments. The principles established by von Frisch, Lorenz, and Tinbergen have been essential for this work.

\subsection{Animal Welfare}

The understanding of animal behaviour has also contributed to animal welfare science. By understanding the behavioural needs of animals, researchers can develop strategies for improving their welfare in captivity and in the wild. The principles established by von Frisch, Lorenz, and Tinbergen have been essential for this work.

\subsection{Inspiration for Future Researchers}

The work of von Frisch, Lorenz, and Tinbergen continues to inspire researchers in ethology, psychology, and neuroscience. Their success in using elegant experiments to answer fundamental questions about animal behaviour demonstrates the power of basic research to improve human health. Their example has motivated generations of scientists to pursue careers in these fields.

\subsection{Recognition and Honors}

In addition to the Nobel Prize, von Frisch, Lorenz, and Tinbergen received numerous other honors and awards for their work. Von Frisch received the Pour le M\'erite in 1965 and was elected to the Royal Society. Lorenz received the Pour le M\'erite in 1964 and was elected to the Royal Society. Tinbergen received the Order of the British Empire in 1973 and was elected to the Royal Society.

Their contributions to ethology and medicine are remembered and celebrated by researchers and clinicians around the world. Their discoveries continue to influence our understanding of animal behaviour and have led to numerous advances in medicine and conservation biology, ensuring that their legacy will endure for generations to come.


\section{Modern Developments}

Frisch, Lorenz, and Tinbergen's work on animal behavior has evolved into modern behavioral ecology and neuroethology. Understanding of innate behavior circuits has informed AI and robotics. Understanding of mate choice has evolutionary implications for human behavior. The study of animal communication has parallels with human language development. Neuroethology---studying the neural basis of natural behavior---has become a major field in neuroscience.
